% (c) 2002 Matthew Boedicker <mboedick@mboedick.org> (original author) http://mboedick.org
% (c) 2003-2007 David J. Grant <davidgrant-at-gmail.com> http://www.davidgrant.ca
% (c) 2008 Nathaniel Johnston <nathaniel@nathanieljohnston.com> http://www.nathanieljohnston.com
% (c) 2011-2026 Scott Clark <scott@scottclark.io> https://scottclark.io
%
%This work is licensed under the Creative Commons Attribution-Noncommercial-Share Alike 2.5 License.

\documentclass[letterpaper,11pt]{article}
\newlength{\outerbordwidth}
\pagestyle{empty}
\raggedbottom
\raggedright
\usepackage[svgnames]{xcolor}
\usepackage{framed}
\usepackage{tocloft}
\usepackage{etoolbox}
\usepackage[hidelinks]{hyperref}
\robustify\cftdotfill

\setlength{\outerbordwidth}{3pt}
\definecolor{shadecolor}{gray}{0.75}
\definecolor{shadecolorB}{gray}{0.93}

\setlength{\evensidemargin}{-0.25in}
\setlength{\headheight}{0in}
\setlength{\headsep}{0in}
\setlength{\oddsidemargin}{-0.25in}
\setlength{\paperheight}{11in}
\setlength{\paperwidth}{8.5in}
\setlength{\tabcolsep}{0in}
\setlength{\textheight}{9.5in}
\setlength{\textwidth}{7in}
\setlength{\topmargin}{-0.3in}
\setlength{\topskip}{0in}
\setlength{\voffset}{0.1in}

\newcommand{\resitem}[1]{\item #1 \vspace{-2pt}}
\newcommand{\resheading}[1]{\vspace{8pt}
  \parbox{\textwidth}{\setlength{\FrameSep}{\outerbordwidth}
    \begin{shaded}
\setlength{\fboxsep}{0pt}\framebox[\textwidth][l]{\setlength{\fboxsep}{4pt}\fcolorbox{shadecolorB}{shadecolorB}{\textbf{\sffamily{\mbox{~}\makebox[6.762in][l]{\large #1} \vphantom{p\^{E}}}}}}
    \end{shaded}
  }\vspace{-5pt}
}
\newcommand{\ressubheading}[4]{
\begin{tabular*}{6.5in}{l@{\cftdotfill{\cftsecdotsep}\extracolsep{\fill}}r}
		\textbf{#1} & #2 \\
		\textit{#3} & \textit{#4} \\
\end{tabular*}\vspace{-6pt}}

\begin{document}

\begin{tabular*}{7in}{l@{\extracolsep{\fill}}r}
\textbf{\Large Scott Clark, Ph.D.} & \textbf{\today} \\
Co-founder \& CEO at Distributional, Inc. & \texttt{scott@scottclark.io} \\
Palo Alto, CA & \texttt{www.scottclark.io} \\
\texttt{github.com/sc932} & \texttt{linkedin.com/in/sc932}
\end{tabular*}
\\

\vspace{8pt}
\noindent \textbf{Note:} This is a longform academic CV, for a shortform resume ($<2$ pages) see ScottClarkResume.pdf \\

\vspace{4pt}
\noindent \textit{Summary.} Three-time founder/executive building ML/AI infrastructure across research, open-source, enterprise, and startup settings. Currently co-founder \& CEO of Distributional (enterprise AI reliability). Previously VP \& GM at Intel (Nov 2020 -- Jun 2023, 200-person AI/HPC Supercomputing org) following Intel's acquisition of SigOpt, where I was co-founder \& CEO (2014--2020). Cornell Ph.D.\ in Applied Mathematics under Prof.\ Peter Frazier. 15 peer-reviewed papers and $\sim$20 granted US patents (1{,}200+ citations, h-index 16), 40+ conference talks and podcast appearances.

%%%%%%%%%%%%%%%%%%%%%%%%%%%%%%
\resheading{Education}
%%%%%%%%%%%%%%%%%%%%%%%%%%%%%%
\begin{itemize}

\item
	\ressubheading{Cornell University}{Ithaca, NY}{Ph.D. Applied Mathematics; M.S. Computer Science}{2008 -- 2012}
	\begin{itemize}
		\resitem{\textit{Department of Energy Computational Science Graduate Fellow} (full 4-year scholarship; DOE contract DE-FG02-97ER25308; $\sim$20 awarded nationally per year).}
		\resitem{Dissertation: ``\textbf{Parallel Machine Learning Algorithms in Bioinformatics and Global Optimization},'' defended Aug 2012.}
		\resitem{Advisor: Prof.~Peter Frazier (ORIE). Committee: Prof.~Steve Strogatz, Prof.~Bart Selman. Research collaborator (not committee): Dr.~Zhong Wang (JGI practicum supervisor + ALE co-author).}
		\resitem{NERSC Startup + Production Allocation as PI (Cray XT4, DOE contract DE-AC02-05CH11231).}
	\end{itemize}

\item
	\ressubheading{Oregon State University}{Corvallis, OR}{B.Sc.~Mathematics, B.Sc.~Computational Physics, B.Sc.~Physics}{2004 -- 2008}
	\begin{itemize}
		\resitem{Graduated \textit{magna cum laude} with three separate Bachelor of Science degrees in four years.}
		\resitem{Minors in Actuarial Sciences and Mathematical Sciences. Paradigms in Physics track.}
		\resitem{Undergraduate thesis: ``Finite Element Modeling on Uncertain Surfaces,'' advisor Prof.~Malgorzata Peszynska (Mathematics); URISC + NSF Grant 0511190. Computational-physics program led by Prof.~Rubin Landau.}
	\end{itemize}

\end{itemize}

%%%%%%%%%%%%%%%%%%%%%%%%%%%%%%
\resheading{Executive \& Industry Experience}
%%%%%%%%%%%%%%%%%%%%%%%%%%%%%%
\begin{itemize}

\item
	\ressubheading{Distributional, Inc. --- \texttt{distributional.com}}{Palo Alto, CA}{Co-founder \& CEO}{July 2023 -- present}
	\begin{itemize}
		\resitem{Analytics for AI agents --- discovering behavioral signals in agent trace data for continuous AI reliability. 30-person team, \$30M raised (Seed Dec 2023 led by Andreessen Horowitz; Series A Oct 2024 led by Two Sigma Ventures).}
		\resitem{Co-founders: Michael McCourt (CTO; multi-paper SigOpt-era co-author), David Rosales (COO), Nick Payton (CRO). 11-person founding team sourced from Bloomberg, Google, Meta, Intel, SigOpt, Slack, Stripe, Uber, Yelp.}
		\resitem{Product evolved 2023--2025 from pre-deployment AI-testing to production behavioral analytics for AI agents --- turning raw trace data into actionable insights for AI teams.}
		\resitem{US Patent 12,505,027 (2025, named inventor): anomaly detection in deployed AI applications.}
	\end{itemize}

\item
	\ressubheading{Intel Corporation}{Portland, OR}{VP \& GM, AI \& HPC Supercomputing Application Level Engineering}{November 2020 -- June 2023}
	\begin{itemize}
		\resitem{Joined Intel through the \$\textit{undisclosed} acquisition of SigOpt (Oct 2020). Led a multi-disciplinary $\sim$200-person organization responsible for application-level AI and HPC software within Intel's Supercomputing Group.}
		\resitem{\textbf{Title progression:} \textit{Director \& GM of SigOpt} (Nov 2020 -- Mar 2021) $\to$ \textit{Senior Director \& GM of SigOpt and AI Application Enablement} (Mar 2021 -- Sep 2022) $\to$ \textit{VP \& GM of AI and HPC Supercomputing Application Level Engineering} (Sep 2022 -- Jun 2023).}
		\resitem{Public face for Intel's oneAPI open-ecosystem thesis --- SC22 theCUBE panel (Dallas, Nov 2022), MLconf 2021 Webinar, Ai4 2021, Intel Conversations in the Cloud Ep.~250.}
		\resitem{Integrated SigOpt's IP into Intel's AI Analytics Toolkit; recognized as a leader in AI Software Optimization (Kisaco Research Leadership Council, 2021).}
	\end{itemize}

\item
	\ressubheading{SigOpt Inc. --- \texttt{sigopt.com}}{San Francisco, CA}{Co-founder \& CEO}{November 2014 -- November 2020}
	\begin{itemize}
		\resitem{Commercial Bayesian-optimization platform serving Fortune 500 enterprises in finance, trading, intelligence, and technology. Acquired by Intel (Oct 29, 2020).}
		\resitem{Raised \$17M across seed + Series A + Series A+ + strategic rounds: Andreessen Horowitz (led seed 2015 and A 2016); Blumberg Capital (lead, 2018 Series A+, with Two Sigma Investments co-investing, announced 2019); Data Collective (DCVC), SV Angel, Stanford, In-Q-Tel strategic, Y Combinator (W15), and notable angels.}
		\resitem{Co-founders: Patrick Hayes (CTO 2014--2021), Eric Liu.}
		\resitem{Grew from Y Combinator W15 through Series A+ with customers in algorithmic trading, government/intelligence, enterprise AI, and consumer-tech.}
		\resitem{Awards: Gartner Cool Vendor in AI Core Technologies (2017); Barclays Open Innovation Challenge winner (2017); Kisaco Research KLC Leader (2021).}
		\resitem{Drove research culture and IP: 15 peer-reviewed papers + $\sim$20 granted US patents as named inventor across Intelligent Optimization Platform, Multi-Criteria Optimization, and Multi-Solution Hyperparameter Tuning families.}
	\end{itemize}

\item
	\ressubheading{Yelp Inc.}{San Francisco, CA}{Data Mining Engineer and Lead on Ad Targeting Team}{July 2012 -- December 2014}
	\begin{itemize}
		\resitem{\textbf{Optimization:} Co-developed and led team for \textbf{MOE} (Metric Optimization Engine, \texttt{github.com/Yelp/MOE}) --- first production Bayesian-optimization open-source package. Third-most-popular Yelp open-source release within 72 hours of launch; basis of SigOpt's founding.}
		\resitem{\textbf{Ad targeting:} Implemented multi-armed bandit strategies for ad selection, sole targeting engineer on mobile ads rollout, developed location-based targeting algorithms.}
		\resitem{\textbf{Director, Yelp Dataset Challenge:} created, implemented, and directed \texttt{yelp.com/dataset\_challenge}; used by 100{,}000+ students worldwide since its inception.}
		\resitem{\textbf{Leadership:} Founded Yelp's internal Applied Learning Group (bi-weekly all-engineering speaker series); MOE team lead; intern and new-hire mentor; 200+ technical interviews.}
		\resitem{Featured in the \textit{Wall Street Journal}, Aug~8 2014 (Elizabeth Dwoskin, ``Big Data's High-Priests of Algorithms'') and the Cornell ORIE alumni spotlight (2014).}
	\end{itemize}

\item
	\ressubheading{Bloomberg LP}{New York, NY}{Financial Software Development Intern}{May -- August 2011}
	\begin{itemize}
		\resitem{Developed end-to-end portfolio-analytics function in C++ and JavaScript (concept $\to$ back-end integration $\to$ GUI $\to$ PDF reporting).}
	\end{itemize}

\end{itemize}

%%%%%%%%%%%%%%%%%%%%%%%%%%%%%%
\resheading{Research Experience (pre-PhD and practicums)}
%%%%%%%%%%%%%%%%%%%%%%%%%%%%%%
\begin{itemize}
\item
	\ressubheading{DOE Joint Genome Institute (LBNL)}{Walnut Creek, CA}{DOE CSGF Practicum under Dr. Zhong Wang \& Rob Egan}{May -- Aug 2010}
	\begin{itemize}
		\resitem{Built open-source genome-assembly validation framework (\textbf{ALE}, published \textit{Bioinformatics} 2013, 208 cites).}
	\end{itemize}
\item
	\ressubheading{Los Alamos National Laboratory}{Los Alamos, NM}{DOE CSGF Practicum under Drs. Nick Hengartner \& Joel Berendzen}{May -- Aug 2009}
	\begin{itemize}
		\resitem{Metagenomics / local sequence-alignment algorithms in Python, C, and CUDA (\texttt{Velvetrope}). This practicum pivoted my dissertation from computational fluid dynamics to bioinformatics.}
	\end{itemize}
\item
	\ressubheading{Max Planck Institute for the Physics of Complex Systems}{Dresden, Germany}{NSF REU under Prof. Steven Tomsovic (WSU)}{May -- Aug 2007}
	\begin{itemize}
		\resitem{Extreme-value statistics of chaotic quantum systems in MATLAB and FORTRAN.}
	\end{itemize}
\item
	\ressubheading{University of California, Davis}{Davis, CA}{NSF REU under Prof. Daniel Cox}{May -- Aug 2006}
	\begin{itemize}
		\resitem{Computational biophysics / protein folding in Java. Results published in \textit{Prion} (2008).}
	\end{itemize}
\end{itemize}

%%%%%%%%%%%%%%%%%%%%%%%%%%%%%%
\resheading{Peer-Reviewed Publications}
%%%%%%%%%%%%%%%%%%%%%%%%%%%%%%
\begin{enumerate}
  \item J Wang, SC Clark, E Liu, PI Frazier. ``Parallel Bayesian Global Optimization of Expensive Functions.'' \textit{Operations Research} 68(6):1850--1865, 2020. \textbf{263 cites.} (arXiv:1602.05149, DBLP journals/ior/WangCLF20.)
  \item I Dewancker, M McCourt, S Clark. ``Bayesian Optimization for Machine Learning: A Practical Guidebook.'' arXiv:1612.04858, 2016. \textbf{142 cites.}
  \item SC Clark. ``Adaptive Sequential Experimentation Techniques for A/B Testing and Model Tuning.'' \textit{WWW 2015 Companion}, p.~911.
  \item I Dewancker, M McCourt, S Clark, P Hayes, A Johnson, G Ke. ``A Stratified Analysis of Bayesian Optimization Methods.'' arXiv:1603.09441, 2016.
  \item I Dewancker, M McCourt, S Clark, P Hayes, A Johnson, G Ke. ``A Strategy for Ranking Optimization Methods Using Multiple Criteria.'' Workshop on Automatic Machine Learning (AutoML, ICML), \textit{PMLR} v.64, 2016.
  \item I Dewancker, M McCourt, S Clark. ``Evaluation System for a Bayesian Optimization Service.'' arXiv:1605.06170, 2016.
  \item SC Clark, R Egan, PI Frazier, Z Wang. ``ALE: a generic assembly likelihood evaluation framework for assessing the accuracy of genome and metagenome assemblies.'' \textit{Bioinformatics} 29(4):435--443, 2013. \textbf{208 cites.}
  \item P Frazier, SC Clark. ``Parallel global optimization using an improved multi-points expected improvement criterion.'' INFORMS Optimization Society Conference, Miami, 2012.
  \item S Clark, E Liu, P Frazier, JL Wang, D Oktay, N Vesdapunt. ``MOE: A global, black-box optimization engine for real-world metric optimization.'' 2014.
  \item KC Kunes, SC Clark, DL Cox, RR Singh. ``Left-handed $\beta$-helix models for mammalian prion fibrils.'' \textit{Prion} 2(2):81--90, 2008.
  \item SC Clark. \textit{Parallel Machine Learning Algorithms in Bioinformatics and Global Optimization.} Ph.D.\ dissertation, Cornell University, Aug 2012.
  \item SC Clark. ``Solving Genomic Jigsaws.'' \textit{DEIXIS Magazine} (DOE CSGF outreach), 2010. (DOE CSGF Communicating Science Award, Honorable Mention.)
  \item \textit{See also:} 5 more SigOpt-era papers and preprints on Bayesian optimization / ranking / MOE (full list available on Google Scholar \texttt{mwWbhAUAAAAJ}).
\end{enumerate}

\textit{Total citations: 1{,}221 (Google Scholar, as of 2026). h-index: 16. i10-index: 28.}

%%%%%%%%%%%%%%%%%%%%%%%%%%%%%%
\resheading{Patents (Selected, as Named Inventor)}
%%%%%%%%%%%%%%%%%%%%%%%%%%%%%%
\noindent $\sim$20 granted US patents across 4 SigOpt patent families (2019--2025) plus 1 Distributional patent (2025). Inventor teams include Patrick Hayes, Michael McCourt, Alexandra Johnson, George Ke, Bolong Cheng, Olivia Kim, Kelvin Tee, Jim Blomo, and others. Selected:
\begin{itemize}
	\resitem{\textbf{Intelligent Optimization Platform} (core SigOpt IP, 11 patents 2019--2025): US 10{,}217{,}061; US 10{,}282{,}237; US 10{,}445{,}150; US 10{,}607{,}159; US 11{,}163{,}615; US 11{,}301{,}781; and others.}
	\resitem{\textbf{Multi-Criteria Optimization} (7 patents 2020--2024): US 10{,}528{,}891; US 10{,}558{,}934; US 10{,}740{,}695; US 12{,}033{,}036; and others.}
	\resitem{\textbf{Multi-Solution Hyperparameter Tuning} (2 patents 2022 \& 2024, plus pending): US 11{,}270{,}217; and others.}
	\resitem{\textbf{Anomaly Detection in Deployed AI Applications} (Distributional, 2025): US 12{,}505{,}027. Inventors: Clark, McCourt, Bourassa-Denis, Laban, Kim, Dewancker, Cheng.}
\end{itemize}

%%%%%%%%%%%%%%%%%%%%%%%%%%%%%%
\resheading{Board \& Civic Service}
%%%%%%%%%%%%%%%%%%%%%%%%%%%%%%
\begin{itemize}
	\resitem{\textbf{Oregon Museum of Science and Industry (OMSI), Board of Trustees} --- Portland, OR, Aug 2022 -- Jul 2025. Treasurer, Executive Committee member, Chair of the Finance Committee.}
	\resitem{\textbf{Oregon State University, College of Science, Board of Advisors} --- May 2019 -- Jun 2025.}
	\resitem{\textbf{Oregon State University, College of Science, Industry and Innovation Council} --- Nov 2019 -- Jun 2025.}
\end{itemize}

%%%%%%%%%%%%%%%%%%%%%%%%%%%%%%
\resheading{Honors, Grants \& Recognition}
%%%%%%%%%%%%%%%%%%%%%%%%%%%%%%
\begin{center}\begin{tabular*}{6.6in}{l@{\extracolsep{\fill}}r}
	\multicolumn{2}{c}{Forbes 30 Under 30, Enterprise Technology \cftdotfill{\cftdotsep}2016}\\
	\multicolumn{2}{c}{OSU College of Science Young Alumni Award \cftdotfill{\cftdotsep}2016}\\
	\multicolumn{2}{c}{Department of Energy Computational Science Graduate Fellowship (DOE CSGF, \$300{,}000) \cftdotfill{\cftdotsep}2008--2012}\\
	\multicolumn{2}{c}{NERSC Production Allocation (PI; 100{,}000 Cray XT4 hours) \cftdotfill{\cftdotsep}2012}\\
	\multicolumn{2}{c}{NERSC Startup Allocation + Renewals (PI; Cray XT4) \cftdotfill{\cftdotsep}2010--2012}\\
	\multicolumn{2}{c}{DOE CSGF Communicating Science Award, Honorable Mention (\textit{DEIXIS} 2010) \cftdotfill{\cftdotsep}2010}\\
	\multicolumn{2}{c}{Cornell University Sage Fellowship (\$55{,}000, declined) \cftdotfill{\cftdotsep}2008--2009}\\
	\multicolumn{2}{c}{Joel Davis Award in Mathematics (OSU, \$1{,}000) \cftdotfill{\cftdotsep}2007--2008}\\
	\multicolumn{2}{c}{URISC Undergraduate Research Fellowship (\$1{,}500) \cftdotfill{\cftdotsep}2007--2008}\\
	\multicolumn{2}{c}{NSF Research Experience for Undergraduates (MPI PKS Dresden) \cftdotfill{\cftdotsep}2007}\\
	\multicolumn{2}{c}{NSF Research Experience for Undergraduates (UC Davis) \cftdotfill{\cftdotsep}2006}\\
	\multicolumn{2}{c}{Paul Copson Memorial Scholarship in Physics (OSU, \$1{,}000) \cftdotfill{\cftdotsep}2006--2007}\\
	\multicolumn{2}{c}{Nicodemus Scholarship in Physics (OSU, \$1{,}000) \cftdotfill{\cftdotsep}2005--2006}\\
	\vphantom{E}
\end{tabular*}\end{center}\vspace*{-16pt}

%%%%%%%%%%%%%%%%%%%%%%%%%%%%%%
\resheading{Selected Invited Talks, Keynotes, \& Conference Appearances}
%%%%%%%%%%%%%%%%%%%%%%%%%%%%%%
40+ documented conference and podcast appearances, 2013--2026. Full catalog with video/audio links and transcripts: \texttt{github.com/sc932}.
\begin{itemize}
	\resitem{\textbf{Distributional / Intel / recent (2021--2026):} SuperComputing SC22 theCUBE Panel on oneAPI (Dallas, Nov 2022); GenAI Week Silicon Valley 2025; AI + a16z Podcast (2025, with Matt Bornstein); CEO.com (2025); BAM ``Coffee with a Founder'' (2025); SAIR Podcast w/ Chuck Ng (Mar 2026); Intel CitC Ep.~250 (2021); MLconf Webinar 2021; Ai4 2021.}
	\resitem{\textbf{SigOpt (2014--2020):} MLconf SF (2014, 2016); MLconf Seattle (2017); NVIDIA GTC Silicon Valley (2017, 2019); O'Reilly AI NY (2019, co-presented with Matt Greenwood, Two Sigma CIO); Ai4 Finance (2019, panel + solo); TWIMLcon (2019); The AI Conference Startup Showcase (2017); Cornell CAM Notable Alumni Speaker Series (Dec 2019, 76-min career-arc talk); SigOpt Summit 2021 keynote; With The Best online conference; ICML 2017 booth.}
	\resitem{\textbf{Yelp / early career (2013--2014):} CMU Silicon Valley TOCS Colloquium (Apr 2013); Yelp Engineering Open House (Nov 2013, ``Optimal Learning for Fun and Profit''); Stanford AI Lab; Microsoft Research; UC Berkeley; Cornell CAM; ACM RecSys 2014 (Palo Alto); Hack Lab Zagreb.}
\end{itemize}

%%%%%%%%%%%%%%%%%%%%%%%%%%%%%%
\resheading{Podcast Guest Appearances (Selected)}
%%%%%%%%%%%%%%%%%%%%%%%%%%%%%%
\begin{itemize}
	\resitem{\textbf{TWIML AI Podcast (Sam Charrington):} Episode 50 (2017), Episode 273 (2019, w/ Matthew Adereth of Two Sigma), Episode 324, TWIMLcon 2019, and 2022 fireside chat. 5+ appearances over 5 years.}
	\resitem{\textbf{Voices in AI with Byron Reese:} Episode 12 (Oct 2017, 56-min long-form).}
	\resitem{\textbf{Emerj / AI in Business Podcast} (Daniel Faggella, 2017); \textbf{IT Visionaries} (2021); \textbf{Intel Conversations in the Cloud} Ep.~250 (2021); \textbf{CEO.com} (2025); \textbf{AI + a16z} w/ Matt Bornstein (2025); \textbf{SAIR} w/ Chuck Ng (Mar 2026).}
\end{itemize}

%%%%%%%%%%%%%%%%%%%%%%%%%%%%%%
\resheading{Press Coverage (Selected)}
%%%%%%%%%%%%%%%%%%%%%%%%%%%%%%
\begin{itemize}
	\resitem{\textit{Wall Street Journal}, ``Big Data's High-Priests of Algorithms / Academic Researchers Find Lucrative Work as `Big Data' Scientists.'' By Elizabeth Dwoskin, Aug 8, 2014.}
	\resitem{Cornell ORIE Alumni Spotlight: ``Scott Clark Ph.D.~'12 brings his work with ORIE Professor Frazier to Yelp and the open source world'' (2014).}
	\resitem{OSU College of Science IMPACT: ``Meet Scott Clark, an OSU science alum who built a \$30M AI startup'' (Tom Henderson, Jul 2025); ``Young Alumni Award Winner Makes Forbes' 30 Under 30 List'' (Nov 2016).}
	\resitem{\textit{DEIXIS Magazine} (DOE CSGF): feature profile (2011, p.~17); author of ``Solving Genomic Jigsaws'' (2010, DOE CSGF Communicating Science Honorable Mention).}
	\resitem{SiliconANGLE / theCUBE coverage of SC22 (Nov 2022, ``VP and GM of AI and HPC Supercomputing Application-Level Engineering'').}
	\resitem{Pulse2 (Amit Chowdhry, Mar 2025) --- Distributional profile interview.}
\end{itemize}

%%%%%%%%%%%%%%%%%%%%%%%%%%%%%%
\resheading{Selected Bylined Writing (beyond peer-reviewed)}
%%%%%%%%%%%%%%%%%%%%%%%%%%%%%%
\begin{itemize}
	\resitem{\textbf{SigOpt Blog} (2015--2020): 23 authored posts on Bayesian optimization and enterprise-AI strategy (\texttt{sigopt.com/blog}, via Wayback).}
	\resitem{\textbf{Yelp Engineering Blog} (2013--2014): 5 authored posts including the canonical MOE launch (``Introducing MOE,'' 2014-07-24) and 3 Yelp Dataset Challenge announcements.}
	\resitem{\textbf{Distributional Blog} (2023--2024): Company fundraise announcements.}
	\resitem{\textbf{NVIDIA Parallel Forall} (2017-08-03, w/ S.~Tartakovsky and M.~McCourt): ``Deep Learning Hyperparameter Optimization with Competing Objectives.''}
	\resitem{\textbf{AWS AI Blog} (2017-05-01, w/ S.~Tartakovsky and M.~McCourt): ``Fast CNN Tuning with AWS GPU Instances and SigOpt.''}
	\resitem{Personal blog (\texttt{scottclark.io/blog}, 2014, Pyramid/Python self-hosted).}
\end{itemize}

%%%%%%%%%%%%%%%%%%%%%%%%%%%%%%
\resheading{Open-Source Projects (\texttt{github.com/sc932})}
%%%%%%%%%%%%%%%%%%%%%%%%%%%%%%
\begin{itemize}
\item \ressubheading{sc932/resume}{LaTeX}{This CV and resume (524+ GitHub stars, briefly \#1 on Hacker News)}{2011 -- present}
\item \ressubheading{Yelp/MOE --- Metric Optimization Engine}{Python, C++, CUDA}{First production Bayesian global-optimization open-source package; basis of SigOpt}{2013 -- 2014}
\item \ressubheading{sc932/ALE --- Assembly Likelihood Estimator}{C, Python}{Probabilistic evaluation of genome assemblies; published \textit{Bioinformatics} 2013 (208 cites); maintained through 2023}{2010 -- 2023}
\item \ressubheading{sc932/Thesis}{LaTeX}{Cornell PhD dissertation source (mirrors Cornell eCommons download)}{2012}
\item \ressubheading{sc932/Velvetrope --- Local Sequence Alignment}{Python, C, CUDA}{LANL 2009 practicum output; Chapter 10 of dissertation}{2009 -- 2010}
\end{itemize}

%%%%%%%%%%%%%%%%%%%%%%%%%%%%%%
\resheading{Areas of Expertise}
%%%%%%%%%%%%%%%%%%%%%%%%%%%%%%
\begin{itemize}
	\resitem{\textbf{Research:} Bayesian optimization, Gaussian processes, optimal learning, multi-armed bandits, hyperparameter tuning, LLM evaluation, agentic-AI systems, AI reliability and testing, experiment design, Monte Carlo methods, bioinformatics and genome assembly, parallel and distributed computing, high-performance computing.}
	\resitem{\textbf{Executive:} P\&L ownership and strategy at 200-person scale; M\&A (sell-side, SigOpt$\to$Intel 2020); fundraising across seed, Series A, and strategic rounds ($\sim$\$47M lifetime); board governance (OMSI Treasurer, OSU College of Science).}
	\resitem{\textbf{Technical:} Python; LLM and agentic-AI systems design, testing, and evaluation; production-ML observability and anomaly detection; scientific and distributed computing at DOE supercomputer scale (NERSC Cray XT4 PI during PhD); data pipelines; MapReduce-scale distributed ML.}
	\resitem{\textbf{Communication:} Several hundred technical talks; extensive podcast circuit; national press quotes; published across industry blogs, peer-reviewed venues, and popular-science outlets.}
\end{itemize}

\end{document}
